%!TEX TS-program = lualatex
%!TEX encoding = UTF-8 Unicode

\documentclass[12pt,addpoints]{exam}

%\printanswers

\usepackage{fontspec}
\def\mainfont{Linux Libertine O}
\setmainfont[Ligatures={TeX, Common}, BoldFont={* Bold}, ItalicFont={* Italic}, Numbers={OldStyle,Proportional}]{\mainfont}
%\setmonofont[Scale=MatchLowercase]{Inconsolata} 
%\setsansfont[Scale=MatchLowercase]{Linux Biolinum O} 
\usepackage{microtype}

\usepackage[left=1in,right=0.75in,top=1in]{geometry}     


\usepackage{graphicx}
	\graphicspath{%
	{/Users/goby/Pictures/teach/348/exams/}}%

% Used to get the last two digits of the year for the header below.
\def\short#1{\csname @gobbletwo\expandafter\endcsname\number#1}

\pagestyle{headandfoot}
\firstpageheader{Marine Biology, Exam 1, \textsc{s}\short{\year}, \numpoints~points.}{}{%
	\ifprintanswers\textbf{KEY}\else Name: \enspace \makebox[2.5in]{\hrulefill}\fi}
\runningheader{}{}{\small{pg. \thepage}}
%\runningheadrule

%\footer{}{}{}%{\makebox[0.75in]{\hrulefill}}
%\extrafootheight{-0.5in}
\footer{\iflastpage{\hspace{-1.2cm}\small \rule{0.6in}{0.4pt} / \numpoints\ points}{}}{}{}%{\makebox[0.75in]{\hrulefill}}
\extrafootheight{-0.5in}

\pointsinmargin
\marginpointname{ pts}
\marginbonuspointname{ pts \textsc{e.c.}}

\usepackage{multicol}

%% Remove comment %% to print answer key.
\unframedsolutions
\renewcommand{\solutiontitle}{}
\SolutionEmphasis{\bfseries}

%% Create a Matching question format
\newcommand*\Matching[1]{
\ifprintanswers
	\textbf{#1}
\else
	\rule{2in}{0.5pt}
\fi
}
\newlength\matchlena
\newlength\matchlenb
\settowidth\matchlena{\rule{2.1in}{0pt}}
\newcommand\MatchQuestion[2]{%
	\setlength\matchlenb{0.98\linewidth}
	\addtolength\matchlenb{-\matchlena}
	\parbox[t]{\matchlena}{\Matching{#1}}\enspace\parbox[t]{\matchlenb}{#2}}


\newcommand*\AnswerBox[2]{%
    \parbox[t][#1]{0.92\textwidth}{%
    \begin{solution}#2\end{solution}}
    \vspace{\stretch{1}}
}

\newenvironment{AnswerPage}[1]
    {\begin{minipage}[t][#1]{0.92\textwidth}%
    \begin{solution}}
    {\end{solution}\end{minipage}
    \vspace{\stretch{1}}}

\newlength{\basespace}
\setlength{\basespace}{5\baselineskip}

%% To hide and show points
\newcommand{\hidepoints}{%
	\pointsinmargin\pointformat{}
}

\newcommand{\showpoints}{%
	\nopointsinmargin\pointformat{(\thepoints)}
}

\newcommand{\bumppoints}[1]{%
	\addtocounter{numpoints}{#1}
}


\begin{document}

\begin{questions}

\vspace*{-1\baselineskip}


\question[2]
Chloride and sodium are the two most common ions in sea water. List the next four most common ions in sea water, from highest to lowest concentration. Do not list the concentration of each ion. 1/2 point per ion.

	\AnswerBox{0.1\textheight}{%
	Sulfate, Magnesium, Calcium, Potassium.
}


\question[2]
How many atmospheres of pressure does a squid living 560 meters deep in the Pacific Ocean experience?

\AnswerBox{\baselineskip}{57 atm}

\question[2]
Name the ocean and the trench or deep within that ocean that is the deepest point on Earth.

\AnswerBox{\baselineskip}{Pacific: Mariana Trench / Challenger Deep}


%\question[4]
%You learned the deepest point of each ocean. Of these five points, list the three \emph{shallowest} points (i.e., least deep), from shallow to deep, and the ocean in which each point is located. 1/2 point per blank plus 1 point for the correct order.
%
%\vspace{0.5\baselineskip}
%
%\begin{tabular}{rll}
%	& Trench Name		&	Ocean Basin \\[4ex]
%	%
%	Most shallow: &
%	\makebox[6cm][l]{\ifprintanswers\textbf{Eurasian Basin / Litke Deep}\else \hrulefill \fi} &
%	\makebox[6cm][l]{\ifprintanswers\textbf{Arctic}\else \hrulefill \fi} \\[4ex]
%	%
%	&
%	\makebox[6cm][l]{\ifprintanswers\textbf{Sandwich}\else\hrulefill \fi} &
%	\makebox[6cm][l]{\ifprintanswers\textbf{Southern Ocean}\else \hrulefill \fi} \\[4ex]
%	%
%	Most deep: &
%	\makebox[6cm][l]{\ifprintanswers\textbf{Sunda/Diamantina}\else \hrulefill \fi} &
%	\makebox[6cm][l]{\ifprintanswers\textbf{Indian}\else \hrulefill \fi} \\
%\end{tabular}

\question[6]
%Briefly explain the difference between planktotrophic and lecithotrophic larval strategies, and why the planktotrophic strategy is less common than the lecithotrophic strategy for marine organisms living at high latitudes. 
Briefly explain the difference between planktotrophic and lecithotrophic larval strategies, and why the lecithotrophic strategy is more common than the planktotrophic strategy for marine organisms living at high latitudes. 

	\AnswerBox{0.3\textheight}{%
	Planktotrophic larvae eat phytoplankton. The phytoplankton season is so short at high latitudes that a slight mistiming of larvae may miss the food supply. Lecithotrophic larvae do not feed so are not dependent on a food source so they have increased chance of survival.
}

\bonusquestion[1]
Tell me what would you \emph{see} if you walked along the abyssal plain from Maine to Spain. Why? This is a trick question; I do have a sense of humor.

\newpage
\vspace*{-1\baselineskip}

\fullwidth{\noindent Use correct names to fill in the blanks below for the zones identified by letters in the figure. (1 point each.)\vspace{-0.5\baselineskip}
\begin{center}\includegraphics[width=0.9\textwidth]{exam1_figure}\end{center}
}%end full width

\question\MatchQuestion{Epipelagic Zone}{A, based on depth.}
\vspace{1\baselineskip}

\question\MatchQuestion{Disphotic Zone}{B, based on light.}
\vspace{1\baselineskip}

\question\MatchQuestion{Bathypelagic Zone}{C, based on depth.}
\vspace{1\baselineskip}

\question\MatchQuestion{Aphotic Zone}{C, D, and E combined, based on light.}
\vspace{1\baselineskip}

\question\MatchQuestion{Neritic Zone}{F, open water out to shelf break.}
\vspace{1\baselineskip}

\question\MatchQuestion{Sublittoral}{G, benthic, from low water mark.}
\vspace{1\baselineskip}

%\question\MatchQuestion{Pelagic Zone}{A and F combined, across ocean from shore to shore.}
%\vspace{1\baselineskip}% Switch to all zones combined, across ocean from shore to shore.

%\question\MatchQuestion{Littoral /Intertidal}{Benthic between low and high water. Not lettered.}
%\vspace{1\baselineskip}

\fullwidth{Extra credit: Proper names for D and E. (1/2 point each.)\vspace*{1\baselineskip}

D. \ifprintanswers \textbf{Abyssopelagic} \else \rule{2.1in}{0.4pt} \fi \qquad \qquad%
E. \ifprintanswers \textbf{Hadopelagic} \else \rule{2.1in}{0.4pt} \fi
}%

%\vspace{1\baselineskip}


%	\ifprintanswers\textbf{Nothing. No light.}\else \vspace*{\baselineskip}\fi}
	

% Points bumped based on how many ocean zone questions used.
% Bump one point for each question.

\bumppoints{6}

\newpage
\vspace*{-1\baselineskip}

\question[10]\label{gladiator}
\begin{multicols}{2}
The gladiator crab \textit{Callinectes gladiator} lives in estuaries along the west coast of Africa (indicated by the dots). The crab larvae leave the estuary to develop in the open ocean. In late spring, the local winds blow from north to south. By later summer, the winds blow from south to north. State which season a crab living in Bengo Bay should release its larvae to ensure the larvae reach the open ocean and then return to the estuary for metamorphosis to adults. Explain why, using the information provided. Include the role of Ekman transport in your answer. % Took out Coriolis %Include the two forces involved and describe why or how they help steer the larvae in the local surface currents.

\columnbreak
	\includegraphics[width=0.47\textwidth]{africa}
\end{multicols}


\begin{AnswerPage}{3\basespace}
	The site is located in the southern hemisphere. 
	\vspace*{0.5\baselineskip}
		
	Therefore, the Coriolis effect causes currents to turn counterclockwise relative to the direction of travel. [2 points for Coriolis and proper direction.]
	\vspace*{0.5\baselineskip}
	
	Ekman transport causes currents to move 90° to the left of the wind direction. [2 points for Ekman and proper direction.]
	\vspace*{0.5\baselineskip}
	
	In late summer, winds blow from south to north, so Ekman transport would cause the currents to move out over the continental shelf. This would be the time for the crabs to release its larvae. [3 points]
	\vspace*{0.5\baselineskip}
	
	In spring, the winds move north to south, so Ekman transport would cause the currents to move towards shore. This is when the larvae will return to the estuary. [3 points]
	
\end{AnswerPage}


\newpage

\question
You may use illustrations to aid your answers to this two-part question. Please do not tell me about earthquakes and volcanos in your answers.

\begin{parts}
	\part[5] Name the two types of the Earth's crust, and then list the differences between them for age, density, and thickness.
	
	\begin{AnswerPage}{0.1\textheight}
	
	Continental crust (1 point) is older (1 point), less dense (1 point), and thicker (1 point) than oceanic crust (1 point).
	
	\end{AnswerPage}

	\part[6] Describe ridge push and how it contributes to the movement of tectonic plates and formation of new sea floor.
	
	\begin{AnswerPage}{0.3\textheight}
	1. Magma rises through the space at the edges of two oceanic crust plates, pushing up the plate edges. [2 points]
	
	\vspace*{0.5\baselineskip}
	
	2. The magma cools on the edge of the plates. [2 points]
	\vspace*{0.5\baselineskip}
	
	3. The weight of the plate edges pushes the plate away from the ridge.  [2 points]
	
	\end{AnswerPage}
	
%	\part[6] Describe slab pull and how it contributes to the movement of tectonic plates and formation of trenches. 
%	
%	\begin{AnswerPage}{0.3\textheight}
%		
%	1. Denser oceanic crust meets and sinks below the less dense continental crust. [2 points]
%	\vspace*{0.5\baselineskip}
%
%	2. As oceanic crust sinks, it becomes even denser. [2 points]
%	\vspace*{0.5\baselineskip}
%
%	3. This causes it to sink faster, pulling along the rest of the plate. [2 points]
%\end{AnswerPage}	
\end{parts}

\question[5]
Briefly describe one common adaptation found among bottom-dwelling organisms living along a passive margin coastline that you are less likely to see in bottom-dwelling organisms living along an active margin. Tell why.

\begin{AnswerPage}{0.3\textheight}
	
\end{AnswerPage}


\newpage
\vspace*{-0.5\baselineskip}

\question
	\begin{parts}
	\part[15]
Summarize the classic model of primary productivity in the North Atlantic Ocean from late winter to early summer.  Your summary should name the major phytoplankton and zooplankton organisms and describe their roles in overall primary production. Your summary should also describe the seasonal availability of nutrients and light to explain the timing of peak productivity.  You may use an illustration to aid your explanation. Do not discuss the fall season.

\begin{AnswerPage}{5\basespace}
1. In the winter, surface temperatures are cold. Winter storms mix surface and deeper water, bringing nutrients to the surface water. [2 points.] 
\vspace*{0.5\baselineskip}

2. As the winter recedes and spring begins, increasing sunlight triggers the bloom. The diatoms and dinoflagellates undergo a bloom, rapidly increasing primary productivity because they now have both nutrients and sufficient light availability. [2 points. Diatoms must be mentioned.]
\vspace*{0.5\baselineskip}

3. During the winter, copepods are down deep. In later spring, the growing adults arrive at the surface, after the bloom and PP has begun. Rapid grazing by copepods (up to 12,000 diatoms per minute) depletes the diatom population, reducing PP. [2 points. Copepods must be mentioned.]
\vspace*{0.5\baselineskip}

4. The nutrients at the surface are eventually depleted by remaining diatoms, shutting down PP to very low levels. As fall approaches, the surface cools again. Early winter storms cause mixing, triggering a second smaller bloom.  This time, the extent of PP is limited by the shorter days as the fall shifts to winter [2 points]
\vspace*{0.5\baselineskip}

[3 points allotted for clarity, proper terminology, grammar and spelling.]
\end{AnswerPage}

	\part[5]
	%Name and describe the modification to the classic model to account for the amount of primary productivity and retention of POM and DOM in tropical surface waters. 
	Name and describe the addition to the classic model to explain how bacteria contribute
	to the retention of \textsc{pom} and \textsc{dom} in surface waters. 

	\begin{AnswerPage}{0.5\basespace}
	
	Microbial loop [1 point]
	\vspace*{0.5\baselineskip}

	Bacteria consume \textsc{dom}. Viruses infect and rupture bacteria, releasing \textsc{om}. [2 points]
	\vspace*{0.5\baselineskip}

	Nanozooplankton consume bacteria, transferring \textsc{om} to higher trophic levels. [2 points]
	
	\end{AnswerPage}
	
	\end{parts}
	

%\question[5]
%Describe the difference(s) between meronektonic and holonektonic organisms. Tell which of these two terms best suits the gladiator crab of question~\ref{gladiator}.
%%Describe the difference(s) between meroplanktonic and holoplanktonic organisms. Tell which of these two terms best suits the gladiator crab of question~\ref{gladiator}.
%
%\begin{AnswerPage}{0.4\textheight}
%	
%\end{AnswerPage}

%\question[5]
%Briefly describe one common adaptation found among bottom-dwelling organisms living along a passive margin coastline that you are less likely to see in bottom-dwelling organisms living along an active margin. 

% NOTE: For a few more points (?) have them describe the difference between passive and active margins.

%\begin{AnswerPage}{0.4\textheight}
%	
%\end{AnswerPage}
%
\newpage

\question
	\begin{parts}
	\part[6]
	On screen is a Pacific Footballfish \textit{(Himantolophus sagamius)}. State whether the footballfish more likely lives in the mesopelagic or bathypelagic zone, and explain one adaptation the fish exhibits to justify your answer. The adaptation \emph{must be characteristic of the depth zone and must be visible in the image.}  The adaptation must be specific to your chosen zone, not common to all deep-sea fishes. Bioluminescence  is \emph{not} visible in the image.
	
	\smallskip
	
	Be sure to explain the specific relationship between the adaptation and the zone you choose.

		\begin{AnswerPage}{2\basespace}
			Bathypelagic [2 points], plus one of:
			\vspace*{0.5\baselineskip}

			\quad Small eyes [2 pts] don't expend energy for lightless zone but still see biolum. [2 points]
			\vspace*{0.5\baselineskip}

			\quad Blobby body {2 pts};  pressure provides support. No energy wasted on skeleton. [2 points]
			\vspace*{0.5\baselineskip}

			\quad Consider color. [2 points plus 2 points for explanation.]
		\end{AnswerPage}

	\part[5]
	Identify a different adaptation that is visible in the image that is found in many deep-sea species, regardless of whether they live in the mesopelagic or bathypelagic. Explain how that adaptation is related to the environment or survival in the deep-sea. Bioluminescence is \emph{not} visible in the image.
	
%	\smallskip
	
%	Be sure to explain how this adaptation might increase survival in the deep sea.

		\begin{AnswerPage}{2\basespace}
			Apparently large stomach to hold large prey. [2 points] \\
			Nutrients are scarce. [3 points]
			\vspace*{0.5\baselineskip}

			Apparently large mouth to capture large prey. [2 points] \\
			Nutrients are scarce. [3 points]
			\vspace*{0.5\baselineskip}

			Consider color. [4 points]
			%Biolum lure to attract prey, which are scarce.
		\end{AnswerPage}

	\end{parts}

\end{questions}

\end{document}